\documentclass[11pt]{article}
\usepackage[margin=1in]{geometry}
\usepackage{amsmath,amssymb,amsthm}
\usepackage{graphicx}
\usepackage[hidelinks]{hyperref}

\newtheorem*{theorem}{Theorem}
\newtheorem*{proposition}{Proposition}
\newtheorem*{remark}{Remark}

\title{A Literature-Implied Negative Answer to a Cuth--Johanis Question}
\author{Automatic Math Discovery Run: FA/Banach 001}
\date{June 10, 2026}

\begin{document}
\maketitle

\section*{Source question}

The source paper is Cuth--Johanis, \emph{Isometric embedding of $\ell_1$ into
Lipschitz-free spaces and $\ell_\infty$ into their duals}, arXiv:1604.04131.
The relevant question appears on page 2:

\begin{center}
\includegraphics[width=\linewidth]{figures/open_question_crop.png}
\end{center}

\section*{Status of this packet}

This packet is filed as a literature-implied answer, not as a new original
pipeline proof. Ostrovskii later proved a necessary and sufficient criterion
for a transportation cost space, equivalently a Lipschitz-free space, to contain
an isometric copy of $\ell_1$; see arXiv:1910.03625. The criterion gives a
short negative answer to Question 2 by one of the uniformly discrete examples
listed in the Cuth--Johanis paper.

\section*{The later criterion}

We use the following theorem of Ostrovskii. A sequence of disjoint pairs
$(x_i,y_i)$ in a metric space has the \emph{minimum matching property} if, for
every $n$, the edges $\{x_i y_i:1\le i\le n\}$ form a minimum-weight perfect
matching on the vertices $\{x_i,y_i:1\le i\le n\}$.

\begin{theorem}[Ostrovskii]
Let $M$ be an infinite metric space. The Lipschitz-free space $\mathcal F(M)$
contains a subspace isometric to $\ell_1$ if and only if $M$ contains a sequence
of disjoint pairs with the minimum matching property.
\end{theorem}

\section*{A uniformly discrete counterexample}

Let $M=\{v_n:n\in\mathbb N\}$ and define, for distinct $i,j$,
\[
  d(v_i,v_j)=1+\frac{1}{\max\{i,j\}}.
\]
This is a metric: if $i<j<k$, then
\[
d(v_i,v_k)\le d(v_i,v_j)+d(v_j,v_k),\quad
d(v_i,v_j)\le d(v_i,v_k)+d(v_j,v_k),
\]
and the remaining triangle inequality is symmetric. Also $d(v_i,v_j)>1$ for
$i\ne j$, so $M$ is uniformly discrete.

\begin{proposition}
The space $\mathcal F(M)$ contains no subspace isometric to $\ell_1$.
\end{proposition}

\begin{proof}
By Ostrovskii's criterion, it is enough to show that no infinite sequence of
disjoint pairs in $M$ has the minimum matching property.

Take any infinite sequence of disjoint pairs. Choose two pairs so that the
indices of the second pair are larger than the indices of the first pair; after
renaming their four vertices, write the indices as
\[
  q_1<q_2<q_3<q_4 .
\]
The two original edges have total weight
\[
  d(v_{q_1},v_{q_2})+d(v_{q_3},v_{q_4})
  =2+\frac1{q_2}+\frac1{q_4}.
\]
The crossed matching has total weight
\[
  d(v_{q_1},v_{q_3})+d(v_{q_2},v_{q_4})
  =2+\frac1{q_3}+\frac1{q_4}.
\]
Since $q_2<q_3$, the crossed matching is strictly lighter. Thus the original
two-edge matching is not a minimum-weight perfect matching on these four
vertices. Consequently no infinite sequence of disjoint pairs in $M$ satisfies
the minimum matching property, and $\mathcal F(M)$ does not contain an
isometric copy of $\ell_1$.
\end{proof}

\section*{Verification code}

The file \texttt{code/verify\_metric\_and\_matching.py} checks, on finite
initial segments, the triangle inequalities and the crossing inequality used in
the proof. These computations are not a substitute for the proof above; they
are included only as a reproducibility check for the chosen metric.

\end{document}
